\documentclass[instructions.tex]{subfiles}
\begin{document}
 
\chapter{Travailler au salut avec crainte}
 
Saint Paul avait bien raison de dire:\og Travaillez à votre salut avec crainte et tremblement.\fg 1.Qui ne craindrait pour une affaire de cette conséquence, dans laquelle il est si difficile de réussir ? 2.Qui ne tremblerait, en voyant que tant de personnes y ont échoué? 

\primo Ce qui doit nous faire craindre, c'est qu'il y a dans la vie une infinité d'écueils et d'occasions capables de nous perdre. Les plus à craindre sont celles dont nous nous défions le moins, et qui viennent de notre propre faiblesse: faiblesse qu'il faut combattre toute la vie; faiblesse si grande, qu'il ne faut qu'une parole d'une personne que nous aimons, un respect humain, une attache aux biens, une amitié qu'on croit innocente, une complaisance ou un regard, pour nous faire tomber et pour nous perdre. 

C'est donc s'aveugler que de vivre sans crainte, comme, s'il était facile de se sauver. La chute des anges, la chutes du premier homme, celle de David, et de Salomon, doivent nous faire comprendre combien il est facile de se perdre. Tant de personnes qui, après avoir vécu dans la sainteté, sont tombées sans se relever, doivent nous apprendre ce que nous avons à craindre. Hélas! que deviendront les roseaux, si les cèdres sont renversés?\og Si le juste sera à peine sauvé\fg, dit l'Écriture,\og que sera-ce du pécheur\fg.

N'est-ce pas cette difficulté du salut et le danger de la damnation qui faisait dire à Job, le plus saint homme qui fût sur terre: Je craignais toutes mes actions\footnote{Verebar omnia opera mea. Job. 6.28}? N'est-ce pas ce qui engageait saint Paul à châtier son corps, crainte d'être réprouvé\footnote{Castigo corpus meum 1.Cor.9.27}? N'est-ce pas ce qui a fait verser tant de larmes et fait pousser des gémissements continuels à saint Jérôme\footnote{Quotidiè Lacrymae, quotidiè gemilus. S.Hier.}? 

Pourquoi tant de jeunes gens se sont-ils retirés des compagnies du monde ? Pourquoi tant de grands seigneurs ont-ils quitté leurs richesses pour mener une vie pénitente dans les monastères ou dans des déserts? Ils vous diront, avec saint Jérôme, que c'est pour ne pas périr dans le monde, que c'est par la crainte de l'enfer, et pour ne pas manquer leur salut\footnote{Ob gehennae metum. S.Hier.}. Il est vrai qu'on peut se sauver sans pratiquer tout ce que les saints ont fait; mais il n'est pas moins vrai qu'on ne se sauve pas sans violence, sans combats, sans crainte, et que l'on ne travaille à son salut qu'autant qu'on craint de se perdre. 

\secundo Si les dangers du salut doivent nous faire craindre, le petit nombre de ceux qui réussissent doit nous faire trembler. Le Sauveur étant interrogé si le nombre de ceux qui se sauvent est petit\footnote{Si pauci sunt qui salvantur. Luc.13.28}, ne fit que cette réponse:\og Faites vos efforts pour entre par la porte étroite, parce que beaucoup de personnes chercheront à y entrer, et ne le pourront pas\fg. Et dans le chapitre septième de saint Matthieu, il ajoute avec une espèce d'étonnement: Oh! que le chemin est étroit qui conduit à la vie! qu'il y en a peu qui le trouvent\footnote{Quam angusta porta et arcta via est, euqe ducit ad vitam, et pauci sunt qui inveniuut eam ! Matth.7.14}!

L'Évangile remarque que le Sauveur prononca ces paroles avec autorité et en maître\footnote{Docens eos, sicut potestatem habens. Matth.7.29}. Il parlait à ses disciples, il parlait à tous, par conséquent aux chrétiens, aux enfans de son Église. Il la compare, cette Église, à un champ rempli de paille et de grain. Il dit que la paille sera séparée, et jetée dans un feu qui ne s´éteindra jamais, et le bon grain sera placé dans le grenier du Père celeste. Or, il y a plus de paille que de bon grain. Le nombre des réprouvés sera donc, même parmi les chrétiens, plus grand que celui des prédestinés.

Peut-on dire, en effet, que le plus grand nombre des chrétiens suit la route du Ciel ? L´Évangile ne dit-il pas que, pour entrer dans le royaume des cieux, il faut se faire violence , aimer le prochain, même ses ennemis, comme soi-même; être humble comme des enfans; qu'il ne faut pas se conformer au siècle; qu'il faut faire pénitence et persévérer jusqu'à la fin; que ceux que Dieu a prédestinés sont conformes à l'image de son fils; que ceux qui lui appartiennent ont crucifié leur chair avec ses passions; que ceux qui ne portent pas leur croix ne sont pas dignes de lui ; qu'il faut détacher son coeur des biens du monde, etc? Plusieurs bons chrétiens, et même plus qu'on ne pense, vivent de la sorte; mais il est clair qu'ils ne font pas le plus grand nombre.

\tertio Puisque la multitude des réprouvés est si grande, pourquoi ne craignez-vous pas d'être de ce nombre ? Et comment peut-il vous échapper une parole, une action, qui ne soient pesées au poids du sanctuaire? Aussitôt que les apôtres eurent ouï cette parole du Sauveur, qu'un d'entre eux devait le trahir, ils furent tous saisis de crainte. Ah! Seigneur, disaient-ils, ne suis-je point ce malheureux\footnote{Numquid ego sum, Domine? Matth.26.22}? Dites la même chose de vous-même, en voyant le grand nombre de ceux qui se perdent.

Si, en marchant sur le bord d'un précipice, vous étiez assuré qu'un de la compagnie dût y périr, ne craindriez-vous pas ce malheur? Avec quelle précaution ne marcheriez-vous pas ? Prenez les mêmes précautions pour votre salut, puisque de tous les hommes qui sont sur la terre, il n'y en aurait qu'un qui dût être réprouvé, vous devriez craindre d'être ce misérable. Toute homme qui ne craint pas ce danger, est déjà perdu ou sur le point de l'être.

Puisque le salut est si difficile et le nombre des réprouvés si grand, qui pourra donc, direz-vous, être sauvé\footnote{Et quis poterit salvus fieri? Matth.19.25}? Rien n'est impossible à Dieu; avec sa grâce vous pouvez tout. Quelque grand que soit le nombre de ceux qui périssent, vous pouvez être un prédestiné, si vous le voulez. Le salut n'est pas plus difficile pour vous que pour tant d'autres qui se sauvent, et qui ont plus d'obstacles que vous.

Si vous craignez efficacement la damnation, votre salut est commencé; si vous persévérez dans cette crainte, votre prédestination est comme assurée\footnote{Salus erit timentibus nomem tuum. Mich.6}. Mais, au contraire, moins vous craindrez, plus il y a à craindre pour vous. On se perd parce qu'on ne craint pas sa perte. 

Voulez-vous vous sauver ? Ouvrez les yeux sur vos dangers, et craignez, non pas d'une crainte servile et naturelle, mais d'une crainte salutaire, qui vous fasse éviter le péché, qui vous porte à Dieu, qui vous fasse craindre de l'offenser et de le perdre. Cette crainte ranimera votre confiance; et plus vous aurez craint de la sorte pendant la vie, plus vous serez assuré de recevoir miséricorde à la mort\footnote{In timore Domini esto tota diem quia habebis spem in novissimo. Prov.23}.

\section{Résolutions}

\primo J'éviterai avec tout le soin possible tout ce qui pourrait être pour moi une occasion de péché.

\secundo Et je combattrai mon penchant dominant sans relâche. 

\tertio Mais quoi qu'il arrive, je ne me laisserai jamais abattre par le découragement, recourant à Dieu au milieu des dangers, m'humiliant profondément devant lui après une faute échappée dans ma faiblesse, en me hâtant d'aller laver mes blessures dans le sacrement institué pour réconcilier les fidèles à Dieu.

\prayer{O mon Dieu! Pénétrez-moi d'une crainte salutaire.}

\end{document}
